\section{Quantum K-theory (Xiaohan Yan and You-Cheng Chou)}%
\label{sec:Quantum K-theory (Xiaohan Yan and You-Cheng Chou)}

\subsection{Overview}%
\label{sub:Overview}

Gromov-Witten invariants are some way of ``counting'' curves in a smooth projective variety \(X\). GW theory produces numbers of the form
\[ \ab<\alpha_1, \ldots, \alpha_n>_{g, n, \beta} \in \Q \]
for various classes \(\alpha_i \in H^*(X)\), genus \(g\), number \(n\) of marked points, and curve class \(\beta \in H_2(X, \Z)\). This has been used to study mirror symmetry, and outside of mirror symmetry has produced birational invariants. Our goal now is to lift this entire story to K-theory. Quantum K-theory produces numbers
\[ \ab<V_1, \ldots, V_n>_{g, n, \beta} \in \Z \]
for \(V_1, \ldots, V_n \in K(X)\). These satisfy various difference equations\footnote{Compare to differential equations in cohomology.} and have found various applications in 3D mirror symmetry and representation theory.

\subsection{K-theory}%
\label{sub:K-theory}

\begin{defn}
    Let \(X\) be a smooth projective variety. Then the \textit{K-theory} \(K(X)\) of \(X\) is the Grothendieck group of vector bundles on \(X\) with the relation that if
    \[ 0 \to V_1 \to V_2 \to V_3 \to 0 \]
    is a short exact sequence, then \(V_2 = V_1 + V_3\).
\end{defn}

\begin{rmk}\leavevmode
    \begin{enumerate}
        \item \(K(X)\) is a ring with the operations of direct sum and tensor product;
        \item The Chern character gives an isomorphism \(K(X) \to A^*(X)\), at least working over \(\Q\).
    \end{enumerate}
\end{rmk}

\begin{exm}
    Let \(X = \P^n\) and set \(P = \mc{O}(-1)\). Then we have
    \[ K(X) = \C[P^{\pm 1}]/(P-1)^{n+1} \xrightarrow{\on{ch}} A^*(X) = \C[H]/H^{n+1}. \]
\end{exm}

For a Deligne-Mumford stack which is not assumed to be smooth, then there are two versions
\[ K_0(X) \supset K^0(X), \]
where \(K_0\) means coherent sheaves and \(K^0\) means perfect complexes. These coincide on a smooth projective variety, so from now on \(K(X)\) will mean \(K_0(X)\). There is a tensor product
\[ K^0(X) \otimes K_0(X) \to K_0(X) \]
and a pushforward
\[ f_* \colon K(X) \to K(Y) \qquad V \mapsto \sum (-1)^i R^i f_* V \]
along proper morphisms. If \(f\) is a perfect morphism (for example flat, regular embedding, etc), then there is a pullback
\[ f^* \colon K(Y) \to K(X) \qquad W \mapsto \sum (-1)^i \on{Tor}_i^{f^{-1} \mc{O}_Y} (\mc{O}_X, f^{-1} W). \]

\subsection{Moduli of curves}%
\label{sub:Moduli of curves}

Recall the moduli space
\[ \bar{M}_{g,n}(X, \beta) = \ab\{ f \colon (\Sigma, p_1, \ldots, p_n) \to X \mid g(\Sigma) = g, f_*[\Sigma] = \beta, \abs{\Aut f} < \infty \} \]
of stable maps to \(X\). This has evaluation maps
\[ \ev_i \colon \bar{M}_{g,n}(X, \beta) \to X \qquad (f, \Sigma, \ul{p}) \mapsto f(p_i). \]
There are also the universal cotangent lines \(L_i = \sigma_i^* \omega_{ \mc{C}/\bar{M} }\).

The moduli space is a proper DM stack which has a perfect obstruction theory, which produces a virtual structure sheaf \(\mc{O}^{\vir}_{g,n,\beta}\). This comes from the structure of a \textit{virtual tangent bundle}, where
\[ \mathbb{T}_{(f, \Sigma, \ul{p})} \bar{M}_{g,n}(X, \beta) = \chi(\Sigma, f^* T_X) - \chi(\Sigma, T_{\Sigma}(-p_1 - \cdots - p_n)). \]
Roughly, if \(\bar{M}\) is the zero locus of a section of a vector bundle \(E\) on a variety \(Y\), then \(\mc{O}^{\vir} = \mc{O}_Y \otimes \mr{Eu}(E)\), where the K-theoretic Euler class is \(\mr{Eu}(E) = \sum (-1)^i \Lambda^i E^{\vee}\).

\subsection{KGW invariants}%
\label{sub:KGW invariants}

For quantum K-theory, see~\cite{Lee2004}; for the permutation-equivariant theory, see~\cite{Givental2017K1,GiventalK7}. To\"en~\cite{Toen1999} proves the Riemann--Roch formula for Deligne--Mumford stacks used below.

\begin{defn}
    K-theoretic GW invariants are defined by
    \[ \ab<V_1, \ldots, V_n>_{g,n,\beta} \coloneqq \chi\ab(\bar{M}_{g,n}(X, \beta); \mc{O}^{\vir} \otimes \bigotimes_{i=1}^n \ev_i^* V_i). \]
\end{defn}

If \(V_1, \ldots, V_n\) are actual vector bundles, then the invariant is an integer, but generally we will consider \(K_0\) with coefficients in a field of characteristic zero. We may also define gravitational descendents
\[ \ab<V_1 L_1^{k_1}, \ldots, V_n L_n^{k_n}>_{g,n,\beta} \coloneqq \chi\ab(\bar{M}_{g,n}(X, \beta); \mc{O}^{\vir} \otimes \bigotimes_{i=1}^n ( \ev_i^* V_i )L_i^{k_i} ). \]
Using the action of \(S_n\) on \(\bar{M}_{g,n}(X, \beta)\) given by permuting the marked points, then the objects \(\bigotimes_{i=1}^n \ev_i^* V\) and \(\bigotimes_{i=1}^n L_i^k\) are \(S_n\)-equivariant.

\begin{defn}
    We define \textit{permutation-invariant K-theoretic GW invariants} by
    \begin{align*}
        \ab<V L_1^k, \ldots, V L_n^k>_{g,n,\beta}^{S_n} &\coloneqq \chi\ab([ \bar{M}_{g,n}(X, \beta)/S_n ]; \mc{O}^{\vir} \otimes \bigotimes_{i=1}^n ( \ev_i^* V ) L_i^k) \\
        &= \dim \chi_{S_n} \ab(\bar{M}_{g,n}(X, \beta); \mc{O}^{\vir} \otimes \bigotimes_{i=1}^n ( \ev_i^* V ) L_i^k )^{S_n}.
    \end{align*}
\end{defn}

Given \(t(q) = \sum_{n \in \Z} t_n q^n \in K(X)[q,q^{-1}]\), we can define
\[ \ab< t(L_1), \ldots, t(L_n)>^{S_n}_{g,n,\beta} \coloneqq \chi ([\bar{M}_{g,n}(X,\beta)/S_n], \cdots). \]
More generally, if \(H = S_{a_1} \times \cdots \times S_{a_m} \subset S_n\), we can define
\[ \ab<t_1(L_1), \ldots, t_1(L_{a_1}), t_2(L_{a_1+1}), \ldots, t_2(L_{a_1+a_2}), \ldots>_{g,n,\beta}^H. \]

\begin{rmk}
    We would like to discuss some relation to ordinary GW invariants. Recall that Hirzebruch-Riemann-Roch states that 
    \[ \chi(M; V) = \int_M \on{ch}(V) \on{td}(T_M). \]
    In the case where \(M\) is a smooth DM stack, T\"oen proved the formula
    \[ \chi(M; V) = \int_{IM} \wt{\on{ch}}(V) \wt{\on{td}}(T_M), \]
    where \(IM = M \times_{\Delta, M \times M, \Delta} M\) is the inertia stack of \(M\) and the tildes mean taking the orbifold Chern character and Todd class. 
\end{rmk}

\begin{exm}
    If \(M = B\mu_r\), then \(IM = \bigsqcup_{g \in \mu_r} B \mu_r\). If \(M = \P(1, r)\) is a weighted \(\P^1\), then \(IM = M \sqcup \bigsqcup_{g \neq 1 \in \mu_r} B\mu_r\).    
\end{exm}

\begin{exm}
    Let \(M = BG\) for a finite abelian group \(G\) and \(V\) be a vector bundle on \(BG\), which is the same thing as a representation of \(G\). Then \(\chi(M, V) = \dim V^G\). On the other hand, we have
    \[ \int_{IM} \wt{\on{ch}}(V) \cdot \wt{\on{td}}(T_M) = \int_{IM} \on{tr}_g V = \sum_{g \in G} \frac{1}{\abs{G}} \on{tr}_g V = \dim V^G, \]
    so we have verified Riemann-Roch in this setting.
\end{exm}

\subsection{Axiomatic properties}%
\label{sub:Axiomatic properties}

\begin{lem}[{\cite{GiventalK7}}]
    We have the \textit{string equation}
    \begin{align*}
        \ab< 1, x_1(L), \ldots, x_k(L), t, \ldots, t >_{g, 1+k+n, \beta}^{S_n} = \ab<x_1(L), \ldots, x_k(L), t^{\otimes n}>^{S_n}_{g,k+n,\beta} \\
        + \sum_i \ab< \cdots, \frac{x_i(L) - x_i(1)}{L-1}, \ldots, t^{\otimes n}>^{S_n}_{g,k+n,\beta}. 
    \end{align*}
\end{lem}

\begin{proof}
    The proof is the same as in GW theory and comes from computing \(\pi^* L_i\) along the map \(\pi \colon\bar{M}_{g,n+1}(X,\beta) \to \bar{M}_{g,n}(X,\beta)\).
\end{proof}

\begin{lem}[{\cite{GiventalK7}}]
    We have the \textit{dilaton equation}
    \begin{align*}
        \ab<L-1, x_1(L), \ldots, x_k(L), t^{\otimes n}>_{g,1+k+n,\beta}^{S_n} ={}& (k-2)\ab<x_1(L), \ldots, x_k(L), t^{\otimes n}>_{g,k+n,\beta}^{S_n} \\
        &+ \ab<x_1(L), \ldots, x_k(L), t, t^{\otimes n-1}>_{g,k+n,\beta}^{S_{n-1}}.
    \end{align*}
\end{lem}

\begin{proof}
    We use the fact that
    \[ \pi_* L = k + \on{Ind}_{S_{n-1}}^{S_n}(1) - 1, \]
    where \(\pi\) forgets the first marked point and \(L = \omega_{\pi} (p_1 + \cdots + p_n)\). This follows from the fact that there is a residue map
    \[ H^0(L) \xrightarrow{\on{Res}} \bigoplus_{i=1}^{k+n} \C p_i \xrightarrow{\sum} \C. \]
    The middle term is a direct sum of \(k\) trivial representations and one standard representation.
\end{proof}

\begin{notn}
    Fix \(\tau, t \in K(X)\). We will denote
    \[ \ab<\!\ab<x_1(L_1), \ldots, x_k(L_k)>\!>_{g,k} \coloneqq \sum_{m,n,\beta} \frac{Q^{\beta}}{m!} \ab<x_1(L_1), \ldots, x_k(L_k), \tau^{\otimes m}, t^{\otimes n}>_{g,k+m+n,\beta}^{S_n}. \]
    We will let \(\{\phi_{\alpha}\}\) be a basis of \(K(X)\) and \(\{\phi^{\alpha}\}\) be the dual basis. Then define the pairing
    \[ g_{\alpha\beta} = \chi(X, \phi_{\alpha} \otimes \phi_{\beta}). \]
    This is not quite the right pairing, so we will consider
    \[ G_{\alpha\beta} = g_{\alpha\beta} + \ab<\!\ab<\alpha,\beta>\!>_{0,2}. \]
    We will let \(g^{\alpha\beta} = (g_{\alpha\beta})^{-1}\) and \(G^{\alpha\beta} = (G_{\alpha\beta})^{-1}\).
\end{notn}

\begin{lem}[{\cite{Givental2000K}}]
    Given \(A,B,C,D \in K(X)[q,q^{-1}]\), we have the \textit{WDVV equation}
    \begin{align*}
        \sum_{\alpha,\beta} \ab<\!\ab< A(L), B(L), \phi_{\alpha}>\! >_{0,3}& G^{\alpha\beta}\ab<\!\ab< C(L), D(L), \phi_{\beta}>\! >_{0,3} \\
        &= \sum_{\alpha,\beta} \ab<\!\ab< A(L), C(L), \phi_{\alpha}>\! >_{0,3} G^{\alpha\beta}\ab<\!\ab< B(L), D(L), \phi_{\beta}>\! >_{0,3}. 
    \end{align*}
\end{lem}

\begin{proof}
    Pull back the equality of the boundary divisors on \(\bar{M}_{0,4} = \P^1\). In K-theory, we need to apply the inclusion-exclusion principle when relating \(\mc{O}_D\) to \(\mc{O}_{C_i}\) when \(D = \bigcup C_i\), which accounts for the appearance of \(G^{\alpha\beta}\).
\end{proof}

\subsection{Loop space formalism}%
\label{sub:Loop space formalism}

\begin{defn}
    The \textit{loop space} \((\mc{K}, \Omega)\) consists of 
    \[ \mc{K} = K(X)(q) \ps{Q}\]
    with the symplectic form
    \[ \Omega(f, g) = - \on{Res}_{q = 0, \infty} \ab<f(q^{-1}), g(q)>_X \frac{\d{q}}{q}. \]
\end{defn}

\begin{rmk}
    There is a polarization \(\mc{K} = \mc{K}_+ \oplus \mc{K}_-\), where
    \[ \mc{K}_+ = K(X)[q, q^{-1}]\ps{Q} \]
    and
    \[ \mc{K}_- = \ab\{ \frac{f}{g} \mid f, g \in K(X)\ps{Q}[q], g(0) \neq 0, \deg f < \deg g\}. \]
    For example, \(\frac{1}{1-q} \in \mc{K}_-\).
\end{rmk}

\begin{defn}
    The \textit{big J-function} \(J \colon \mc{K}_+ \to \mc{K}\) is given by
    \[ \mathbf{t} = \mathbf{t}(q) \mapsto 1 - q + \mathbf{t}(q) + \sum_{\beta, \alpha} Q^{\beta} \phi^{\alpha} \ab< \frac{\phi_{\alpha}}{1-qL}, \mathbf{t}(L)^{\otimes n}>_{0, 1+n, \beta}^{S_n}. \]
\end{defn}

\begin{rmk}
    Note that applying Riemann-Roch, we have
    \begin{align*}
        \chi\ab([\bar{M}/S_n], \frac{1}{1-qL}) &= \int_{I[\bar{M}/S_n]} \on{ch}\ab(\frac{1}{1-qL}) \on{td}.
    \end{align*}
    Each term becomes of the form \(\on{ch}\ab(\on{tr}_g \frac{1}{1-qL})\), which gives us something of the form \(\on{ch}\ab(\frac{1}{1-q\zeta L})\) for some root of unity \(\zeta\). Then we can expand
    \[ \frac{1}{1-q\zeta L} = \frac{1}{(1-q \zeta) - q\zeta (L-1)} = \sum_k \frac{(q\zeta)^k}{(1-q\zeta)^{k+1}}(L-1)^k. \]
\end{rmk}

\begin{defn}
    The \textit{Givental cone} is defined to be \(\mc{L} \coloneqq \Ima J \subset \mc{K}\), where we are considering the germ near \(1-q\) in the \(Q\)-adic topology.
\end{defn}

\begin{exm}
    Let \(X = \P^n\). Then we see that
    \[ J(0) = (1-q) \sum_{d \geq 0} \frac{Q^d}{\prod_{\ell=1}^d (1-Pq^{\ell})^{n+1}}, \]
    where \(P = \mc{O}(-1)\). As before, we can expand
    \[ \frac{1}{1-P q^{\ell}} = \frac{1}{(1-q^{\ell}) - q^{\ell}(P-1)} = \sum_{k=0}^{\infty} \frac{q^{\ell k}}{(1-q^{\ell})^{k+1}}(P-1)^k. \]
\end{exm}

\begin{prop}[{\cite{GiventalK7}}]
    In fact, \(\mc{L}\) is actually a cone in \(\mc{K}\) and we have
    \[ \mc{L} = \bigcup_{t \in K(X)\ps{Q}} (1-q) S_{t}^{-1}(q) \mc{K}_+. \]
\end{prop}

We may also introduce a ``huge'' J-function
\[ J(\mathbf{t}, \tau) = 1-q + \mathbf{t}(q) + \tau + \sum_{\beta, k, \alpha} \frac{Q^{\beta}}{k!} \phi^{\alpha} \ab< \frac{\phi_{\alpha}}{1-qL}, \mathbf{t}(L)^{\otimes n}, \tau^{\otimes k}>_{0,1+n+k,\beta}^{S_n}. \]
If we fix \(\mathbf{t}(q)\), then \(J(\tau)\) comes from a fundamental solution of a \textit{differential equation} on the Frobenius manifold in \(\tau\). In addition, \(\mc{L}^{\mathbf{t}}=  \Ima J(\tau)\) is actually a Lagrangian cone.

\begin{exm}
    Consider \(f \colon \C[\Lambda_1^{\pm}, \Lambda_2^{\pm}] \to \C[\Lambda_1^{\pm}, \Lambda_2^{\pm}]\) which is given by the Adams operation
    \[ t \mapsto \on{tr}_{(1\cdots n)} t^{\otimes n} = \Psi^n (t). \]
    Therefore, this sends \(\Lambda_1 \mapsto \Lambda_1^n\) and \(\Lambda_2 \mapsto \Lambda_2^n\). However, it is not clear how to take derivatives in \(t\) because we would need \(nv t^{n-1} = \Psi^n(v)\), so the J-function does not satisfy a quantum differential equation in \(\mathbf{t}\).
\end{exm}

\subsection{Equivariant K-theory and localization}%
\label{sub:Equivariant K-theory and localization}

For the fixed-point recursion and the projective-space calculation, see~\cite{GiventalK2}.

Let a torus \(T = (\C^{\times})^r\) act on a smooth projective variety \(X\).

\begin{defn}
    The \textit{equivariant K-theory} \(K_T(X)\) is the Grothendieck group of \(T\)-equivariant vector bundles on \(X\).
\end{defn}

\begin{exm}
    For example, we have
    \[ K_T(\on{pt}) = \on{Rep}(T) = \C[\Lambda_1^{\pm}, \ldots, \Lambda_r^{\pm}], \]
    where the \(\Lambda_i\) are the \(1\)-dimensional representations with weight \((0, \ldots, 1, \ldots, 0)\).
\end{exm}

\begin{rmk}
    For any \(X\), \(K_T(X)\) is a \(K_T(\mr{pt})\)-algebra.
\end{rmk}

\begin{exm}
    Consider the fully equivariant K-theory of \(\P^n\). Then we have
    \[ K_T(\P^n) = \C[P^{\pm}, \Lambda_0^{\pm}, \ldots, \Lambda_n^{\pm}] \big/ \prod_{i=0}^n (P-\Lambda_i). \]
    When we specialize \(\Lambda_i = 1\), we recover the non-equivariant K-theory.
\end{exm}

\begin{defn}
    We can define the equivariant Euler characteristic as the pushforward
    \[ \chi_T = \pi_* \colon K_T(X) \to K_T(\mr{pt}). \]
\end{defn}

\begin{exm}
    For example, we have \(\chi_T(\mc{O}_{\P^1}(1)) = \Lambda_0^{-1} + \Lambda_1^{-1}\). Similarly, we have
    \[ \chi_T(\mc{O}(2)) = \Lambda_0^{-2} + \Lambda_0^{-1} \Lambda_1^{-1} + \Lambda_1^{-2}. \]
\end{exm}

\begin{notn}
    Note that \(\on{Frac}(T) = \on{Frac}(\on{Rep}(T))\). Then we will set
    \[ K_T^{\loc}(X) \coloneqq \on{Frac}(T) \otimes_{\on{Rep}(T)} K_T(X). \]
\end{notn}

We will need the following fact. Let \(F\) be the fixed point set of \(T\) on \(X\) and consider \(\iota \colon F \hookrightarrow X\). Then there is an isomorphism of \(\on{Frac}(T)\)-algebras
\[ K_T^{\loc}(X) \xrightarrow{\iota^*} K_T^{\loc}(F). \]

\begin{rmk}
    The inverse is given by
    \[ (\iota^*)^{-1} = \bigoplus_{i=1}^r \frac{\iota_{i*}}{\on{Eu}(N_{F_i/X})}, \]
    where \(\on{Eu}(V) \coloneqq \sum_{j = 0}^{\on{rk} V} (-1)^j \Lambda^j V^*\) is the K-theoretic Euler class.
\end{rmk}

\begin{exm}
    On \(\P^n\), we have \(n+1\) isolated fixed points \(F_i\), and then we obtain
    \[ \frac{\iota_{i*} \mc{O}_{F_i}}{\on{Eu}(N_{F_i/\P^n})} = \prod_{j \neq i} \frac{P - \Lambda_j}{\Lambda_i - \Lambda_j}. \]
    These are idempotents of \(K_T^{\loc}(\P^n)\).
\end{exm}

Using equivariant localization, we obtain
\[ \chi_T(X; V) = \sum_i \chi\ab(F_i; \frac{\iota_i^* V}{\on{Eu}(N_{F_i/X})}). \]

\begin{rmk}
    In the \(\Lambda = 1\) limit, we can recover the non-equivariant Euler characteristic.
\end{rmk}

\begin{exm}
    Let \(X = \P^1\). Then we compute
    \begin{align*}
        \chi_T(\P^1, \mc{O}(1)) &= \chi\ab(F_0, \frac{\mc{O}(1)|_{F_0}}{\on{Eu}(T_{F_0/\P^1})}) + \chi\ab(F_1, \frac{\mc{O}(1)|_{F_1}}{\on{Eu}(T_{F_1/\P^1})}) \\
        &= \frac{\Lambda_0^{-1}}{1 - \frac{\Lambda_0}{\Lambda_1}} + \frac{\Lambda_1^{-1}}{1 - \frac{\Lambda_1}{\Lambda_0}} \\
        &= \Lambda_0^{-1} + \Lambda_1^{-1}. 
    \end{align*}
\end{exm}

From now on, we will assume that \(X\) is a GKM manifold. Our goal now is to use localization on \(\bar{M}_{g,n}(X, \beta)\)  to study the big J-function of \(X\). Any \(T\)-fixed stable map \(f \colon (\Sigma, \ul{p}) \to X\) has two types of irreducible components. There are those components which map entirely to fixed points and those components who map via ramified covers of \(1\)-dimensional orbits. Unfortunately, we cannot compute the big J-function via virtual localization, but we can still study it.

Let \(f = f(q) \in \mc{K} = K_T^{\loc}(X)(q)\ps{Q}\). Denote the fixed points by \(F_{\alpha}\) for \(\alpha = 1, \ldots, r\) with inclusions \(\iota_{\alpha} \colon F_{\alpha} \hookrightarrow X\). Then denote
\[ f_{\alpha}(q) = \iota_{\alpha}^* f(q) \in K_T^{\loc}(\on{pt})(q)\ps{Q}. \]
Then \(f \in \mc{L}\) if and only if for all \(\alpha\), the following conditions hold:
\begin{enumerate}
    \item Given \(\lambda \in T_{F_{\alpha}} X\) and \(m \in \N_+\), \(f_{\alpha}(q)\) has at most simple poles at \(q = \lambda^{1/m}\), and
        \[ \on{Res}_{q = \lambda^{1/m}} f_{\alpha}(q) \frac{\d{q}}{q} = -\frac{Q^m}{m} \cdot \frac{\on{Eu}(T_{F_{\alpha}/X})}{\on{Eu}(T^{\vir}_{\phi_{\alpha,\alpha',m}}\bar{M}_{0,2}(X, m \cdot \beta(\alpha, \alpha')))} \cdot f_{\alpha'}(\lambda^{1/m}). \]
        Here, \(\lambda \in T_{F_{\alpha}}X\) means that \(\lambda\) is a \(1\)-dimensional subrepresentation of \(T_{F_{\alpha}} X\). Moreover, \(\beta(\alpha,\alpha')\) is the degree of the \(1\)-dimensional orbit of \(T\) determined by the weight \(\lambda\) and \(\phi_{\alpha, \alpha', m}\) parameterizes \(m:1\) covers of this \(1\)-dimensional orbit.
    \item \(f_{\alpha}(q)\) has its other poles at \(0, \infty\), and roots of unity, and is regular elsewhere. In addition, there are various minor conditions.
\end{enumerate}

\begin{exm}
    Let \(X = \P^n\). Then
    \[ J_T^X(0) = (1-q) \sum_{d \geq 0} \frac{Q^d}{\prod_{i=0}^n \prod_{\ell = 1}^d \ab(1 - \frac{P}{\Lambda_i} q^{\ell})} \in K_T(X)(q) \ps{Q}. \]
    This then implies that 
    \[ J_T^X(0) |_{F_j} = (1-q) \sum_{d \geq 0} \frac{Q^d}{\prod_{i=0}^n \prod_{\ell = 1}^d \ab(1 - \frac{\Lambda_j}{\Lambda_i}q^{\ell})}, \]
    so the poles are at \(q = \ab(\frac{\Lambda_i}{\Lambda_j})^{1/\ell}\). Note also that
    \[ T_{F_j} X = \frac{\Lambda_0}{\Lambda_j} + \cdots + \frac{\Lambda_n}{\Lambda_j} - 1. \]
\end{exm}

\subsection{Quantum K/GV correspondence}%
\label{sub:Quantum K/GV correspondence}

The quintic correspondence was conjectured in~\cite{GaroufalidisScheidegger2022} and proved in~\cite{ChouLeeQuintic}. The following Calabi--Yau threefold formula is due to Chou--Lee~\cite{ChouLee2025}.

Our goal here will be to compute \(J^K(0)\) on a Calabi-Yau threefold. The following result was conjectured by Jockers-Mayr and by Garoufalidis-Scheidegger.

\begin{thm}[{\cite{ChouLee2025}}]
    We have
    \begin{align*}
        J^K(0) &\coloneqq (1-q) + \sum_{\beta, \alpha} Q^{\beta} \Phi^{\alpha} \ab<\frac{\Phi_{\alpha}}{1-qL}>_{0,1,\beta}^{X,K} \\
        &= (1-q) + \sum_{\beta} \sum_{r=1}^{\infty} \mr{GV}_{\beta} Q^{r\beta} \ab[\sum_{j=1}^{n_1} \Phi^{1j} a(r, q^r) + \Phi^{01} b(r, q^r)],
    \end{align*}
    where \(\ab\{\Phi_{\alpha}\}_{\alpha=1}^N\) is a basis of \(K^0(X)\) decomposed such that
    \[ \ab\{\Phi_{\alpha}\}_{\alpha=1}^N = \bigsqcup_{i=0}^3 \ab\{\Phi_{ij}\}_{j=1}^{n_i} \]
    such that the collection \(\ch(\Phi_{ij})\) for fixed \(i\) form a basis of \(H^{2i}(X)\). We then select a dual basis \(\ab\{\Phi^{\alpha}\}_{\alpha} = \ab\{\Phi^{ij}\}_{i,j}\). Moreover, we set the multiple cover formulae
    \begin{align*}
        a(r, q^r) &= \frac{r-1}{1-q^r} + \frac{1}{(1-q^r)^2}, \\
        b(r, q^r) &= \frac{r^2-1}{1-q^r} + \frac{3}{(1-q^r)^2} - \frac{2}{(1-q^r)^3}.
    \end{align*}
    Finally, GV denotes Gopakumar-Vafa invariants, which are defined in genus zero by
    \[ \mr{GW}_{0,\beta} = \sum_{r \mid \beta} \frac{1}{r^3} \mr{GV}_{\beta/r}. \]
    By M\"obius inversion, we see that
    \[ \mr{GV}_{\beta} = \sum_{r \mid \beta} \frac{\mu(r)}{r^3} \mr{GW}_{0,\beta/r}. \]
\end{thm}

\begin{cor}[{\cite{ChouLee2025}}]
    We have
    \[ \ab<1>_{0,1,\beta}^K = \sum_{r \mid \beta} r^2 \mr{GV}_{\beta/r}. \]
    Applying M\"obius inversion, we obtain\footnote{Note that this result was obtained in all genus by Ionel--Parker~\cite{IonelParker2018}.}
    \[ \mr{GV}_{\beta} = \sum_{r \mid \beta} r^2 \mu(r) \ab<1>_{0,1,\beta/r}^K \in \Z. \]
    In addition, we can recover the usual GW invariants by\footnote{This is like if we apply Hirzebruch-Riemann-Roch to compute the K-theoretic pairing, the top degree term only sees the Chern characters of the two inputs.}
    \[ \lim_{q \to 1} (1-q)^3 \ab<\frac{1}{1-qL}>_{0,1,\beta}^{X,K} = \sum_{r \mid \beta} \frac{-2}{r^3} \mr{GV}_{\beta/r} = -2 \mr{GW}_{\beta} = \ab<\psi>_{0,1,\beta}^{\mr{GW}}. \]
\end{cor}

Naively, we can think of \(J^K(0) \in \Z\ps{q}\ps{Q}\), but to be a bit more precise, we can consider
\[ J^K(0) \in 1-q + \mc{K}_-. \]
Here, we expand \(\frac{1}{1-qL} = \sum q^i L^i\). 
We also need to consider the fact that if \(M\) is a DM stack and \(L\) is a line bundle on \(M\), then the function
\[ f^{\Phi}(q) \coloneqq \chi\ab(M, \frac{\Phi}{1-qL}) \in \Z \ps{q} \]
converges to a rational function for \(\abs{q} < \ep\). Note that there are poles only at roots of unity.

This fact can be proved using the orbifold Riemann-Roch theorem. Decompose the inertia as
\begin{align*}
    IM &= \bigsqcup_{i \in I} M_i \\
    &= \bigsqcup_{\zeta} M_{\zeta},
\end{align*}
where \(M_{\zeta}\) is the union of twisted sectors such that \(g\) acts on \(L\) with eigenvalue \(\zeta\). Using this, we can compute
\begin{align*}
    \chi\ab(M, \frac{1}{1-qL}) &= \sum_{i \in I} \int_{M_i} \on{td}(M_i) \on{ch} \ab(\frac{ \on{Tr}_g \ab(\frac{1}{1-qL}) }{\lambda_{-1}(N_{M_i/M}^*)}) \\
    &= \sum_{\zeta} \int_{M_{\zeta}} \on{td}(M_{\zeta}) \sum_{k \geq 0} \frac{(\zeta q)^k}{(1-\zeta q)^{k+1}} \frac{( \on{ch}(L)-1 )^k}{\ch \lambda_{-1}(N^*)},
\end{align*}
and therefore all poles must be at roots of unity. Here, we use the expansion
\[ \frac{1}{1-\zeta q L} = \sum_{k \geq 0} \frac{(\zeta q)^k}{(1-\zeta q)^{k+1}} (L-1)^k \]
and take the Chern character.

\begin{exm}
    Let \(L = \mc{O}(1)\) on \(\P^1\). Then we have
    \begin{align*}
        \chi\ab(\P^1, \frac{1}{1-qL}) &= \chi\ab(\P^1, \frac{1}{1-q} + \frac{q}{(1-q)^2}(L-1)) \\
        &= \sum_{i \geq 0} q^i (i+1) \\
        &= \frac{1}{(1-q)^2}.
    \end{align*}
\end{exm}

\begin{exm}
    We will now consider weighted projective space. Then we have
    \begin{align*}
        \chi\ab(\P(w_0, \ldots, w_n), \frac{1}{1-qL}) &= \prod_{i=0}^n \frac{1}{1-q^{w_i}}. 
    \end{align*}
\end{exm}

